\chapter{网络外部性与需求方规模经济}

\begin{examguide}
\textbf{本章考情}：网络外部性是数字经济微观机制的第一块基石，0258 各校专业课必考（\important{专硕·核心}），学硕微观命题中作为"外部性与市场失灵"的新形式出现（\important{学硕·热点}）。命题形式以推导题与简答题为主：给出效用函数 $u = \theta n - p$ 推导需求函数与临界规模，是计算题的标准素材。

\textbf{核心考点}：（1）直接与间接网络外部性的区分；（2）需求方规模经济与供给方规模经济的对照；（3）$u=\theta n - p$ 模型的需求推导、临界规模与多重均衡；（4）Katz–Shapiro 预期满足均衡；（5）路径依赖与转换成本锁定：迁移壁垒的构成与绝对锁定条件；（6）兼容性决策；（7）网络垄断与传统垄断的成因与后果对照。

\textbf{本章主线}：什么是网络外部性（3.1）→ 如何建模（3.2）→ 兼容与锁定的决策（3.3）→ 对市场结构意味着什么（3.4）。
\end{examguide}

第 2 章准备的最优化与均衡分析工具，从本章起进入实质应用。网络外部性是数字经济区别于传统经济的第一个微观机制：传统商品的价值由其自身属性决定，网络产品的价值却取决于有多少人在用。这一差别改变了需求曲线的形状与均衡的个数，也决定了数字市场为何容易走向集中。

\section{网络外部性的概念}

\subsection{从外部性到网络外部性}

外部性（externality）是经济学的经典概念：当一个经济主体的行为直接影响了其他主体的效用或生产可能性，而这一影响未经过市场价格机制传导时，便存在外部性。马歇尔（Marshall）最早讨论产业集聚带来的"外部经济"（external economies），庇古（Pigou）则将其系统化为市场失灵理论的核心范畴：正外部性导致供给不足，负外部性导致供给过度。

\begin{definition}[网络外部性（network externality）]
某一产品的价值随使用该产品的用户数量的增加而增加，此时每个用户的加入都直接提高了存量用户的效用。
\end{definition}

网络外部性是外部性概念在"网络产品"上的具体形态，电话网络、操作系统、社交平台皆属此类。与传统外部性（污染、灯塔）相比，网络外部性有三个显著差异。其一，它是\textbf{消费侧}的外部性：作用于需求方，改变的是消费者的支付意愿而非厂商的成本；其二，它\textbf{可经由市场机制内部化}：厂商可以通过定价策略（补贴、会员制）部分地将外部性内部化，这正是平台定价理论（第 4 章）的出发点；其三，它天然具有\textbf{自我强化}（self-reinforcing）性质：用户增加提升价值，价值提升吸引更多用户，形成正反馈循环。

\begin{misconception}
\textbf{误区}：把"网络效应"与"网络外部性"当作同义词。严格地说，网络效应（network effects）是用户规模影响产品价值的客观现象；网络外部性则要求这一影响\textbf{未通过市场机制内部化}。若网络所有者能通过定价将规模收益收入囊中（如平台向两侧收费），网络效应存在但并不构成外部性——这正是"可经由市场机制内部化"的含义。多数教材不严格区分二者，但这一区分在学理上有意义。
\end{misconception}

\subsection{直接网络外部性与间接网络外部性}

按外部性的传导路径，网络外部性分为两类：直接网络外部性经用户之间的直接互动传导，间接网络外部性经互补品市场传导。二者的区分直接影响后文的建模方式与厂商策略。

\begin{definition}[直接网络外部性（direct network externality）]
某一产品的价值直接取决于使用该产品的用户数量，即用户效用函数中直接包含用户规模。
\end{definition}

电话网络、微信、传真机是直接网络外部性的典型。对使用微信的消费者而言，其效用直接取决于微信的用户规模：好友越多，微信对他越有用。直接网络外部性的来源是\textbf{用户之间的直接互动}：网络的本质功能就是连接用户，连接对象越多，网络对每个成员的价值越高。

\begin{definition}[间接网络外部性（indirect network externality）]
产品的价值通过互补品的规模间接取决于用户数量，即一侧用户规模的增加通过降低互补品价格或提高互补品质量，间接提升另一侧用户的效用。
\end{definition}

间接网络外部性中，用户之间并不直接互动，其受益路径经过一个市场中介。游戏主机是标准例子：玩家数量增加吸引更多游戏开发商为主机开发游戏，游戏数量增加又提升主机的价值：玩家与开发商互不相识，但通过"主机—游戏"的互补关系相互创造价值。\mn{考试辨析高频：间接网络外部性必须经由互补品市场传导；直接网络外部性来自用户间直接互动。判断方法：删除"互补品市场"这一环节后，外部性是否仍然存在？}间接网络外部性是第 4 章双边市场理论的直接前身：当互补品市场与产品市场的相互作用足够强时，两个市场合并分析即为双边市场模型。

\begin{misconception}
\textbf{误区}：将"间接网络外部性"等同于"交叉网络外部性"。二者相关但层次不同：间接网络外部性描述的是\textbf{产品}层面的现象（互补品规模提升产品价值），交叉网络外部性（cross-group network externality）是双边市场模型中\textbf{两侧用户群体}之间的外部性（A 侧用户规模影响 B 侧效用）。交叉网络外部性是间接网络外部性在平台场景下的精确化，但并非所有间接网络外部性都构成交叉网络外部性。
\end{misconception}

网络外部性并非必然为正。拥堵是典型的负网络外部性（negative network externality）：高速公路上车辆增加使每辆车速度下降。数字场景中同样存在：网络平台上的信息过载、垃圾信息使大网络的用户体验劣于精心运营的小社区。术语使用上，"网络外部性"通常默认指正网络外部性。

\subsection{需求方规模经济与供给方规模经济}

传统规模经济（economies of scale）作用于\textbf{供给方}：产量扩大摊薄固定成本，平均成本曲线下降。其来源是生产技术与成本结构（设备、研发的不可分性）。网络外部性则创造需求方的规模经济。

\begin{definition}[需求方规模经济（demand-side economies of scale）]
需求方规模经济是指产品价值随用户规模扩大而上升，使需求曲线本身向外移动的规模经济形态，其来源在消费侧（网络外部性），与来自生产侧（成本摊薄）的供给方规模经济（supply-side economies of scale）相对。
\end{definition}

两种规模经济可以分别用一个公式刻画。\mn{辨析口诀：供给方规模经济看"成本曲线怎么降"（生产侧），需求方规模经济看"需求曲线怎么移"（消费侧）。}供给方规模经济体现在成本函数上：设固定成本为 $F$、边际成本为常数 $c$，产量 $Q$ 处的平均成本为
\[
AC(Q) = \frac{F}{Q} + c, \qquad AC'(Q) = -\frac{F}{Q^2} < 0,
\]
产量越大，固定成本摊得越薄，平均成本越低并逐渐趋近边际成本 $c$，这正是"成本曲线怎么降"。

需求方规模经济体现在需求函数上：沿用后文 3.2 节基本模型的设定（推导 3.1），给定预期规模 $n^e$，需求为 $n = 1 - p/n^e$，对 $n^e$ 求偏导得
\[
\frac{\partial n}{\partial n^e} = \frac{p}{(n^e)^2} > 0,
\]
预期用户规模越大，同一价格 $p$ 下的需求量越大。把同一关系整理成反需求形式看得更清楚：
\[
p = n^e (1-n),
\]
给定用户规模 $n$，消费者愿意支付的价格与预期规模 $n^e$ 成正比，$n^e$ 上升使整条需求曲线向外移动，这正是"需求曲线怎么移"。图~\ref{fig:eos} 给出二者的直观对照：左图的平均成本曲线随产量下行并趋近边际成本，右图的需求曲线随预期规模上升而整体外移。

\begin{figure}[htbp]
\centering
\begin{tikzpicture}
\begin{axis}[
    width=5.8cm, height=4.2cm,
    xlabel={产量 $Q$}, ylabel={平均成本 $AC$},
    xmin=0, xmax=1.05, ymin=0, ymax=0.32,
    xtick={0,0.5,1}, ytick={0,0.1,0.2,0.3},
    grid=major, axis lines=left,
]
\addplot[color=main, very thick, domain=0.06:1.05, samples=100] {0.012/x + 0.08};
\addplot[color=second, dashed, thick, domain=0:1.05] {0.08};
\node[anchor=north, font=\small, color=second] at (axis cs:0.95,0.077) {$c$};
\addplot[color=winered, dashed, thick] coordinates {(0.5,0) (0.5,0.32)};
\node[anchor=south, font=\small, color=winered, fill=white, inner sep=2pt] at (axis cs:0.5,0.25) {$Q_{\text{MES}}$};
\end{axis}
\end{tikzpicture}
\hfill
\begin{tikzpicture}
\begin{axis}[
    width=5.8cm, height=4.2cm,
    xlabel={用户规模 $n$}, ylabel={价格 $p$},
    xmin=0, xmax=1.05, ymin=0, ymax=0.30,
    xtick={0,0.5,1}, ytick={0,0.1,0.2},
    grid=major, axis lines=left,
]
\addplot[color=main, very thick, domain=0:1, samples=2] {0.15*(1-x)};
\addplot[color=winered, very thick, domain=0:1, samples=2] {0.25*(1-x)};
\draw[->, color=second, thick] (axis cs:0.3,0.113) -- (axis cs:0.3,0.167);
\node[anchor=south, font=\small, color=second] at (axis cs:0.3,0.22) {$n^e$ 上升};
\node[anchor=north, font=\small, color=main] at (axis cs:0.55,0.062) {$n^e_1$};
\node[anchor=south, font=\small, color=winered] at (axis cs:0.72,0.078) {$n^e_2$};
\end{axis}
\end{tikzpicture}
\caption{供给方规模经济与需求方规模经济的对照}
\label{fig:eos}
\end{figure}

二者的区分带来截然不同的市场结构含义。供给方规模经济的优势在市场足够大时会衰减殆尽。\textbf{最低有效规模}（minimum efficient scale，MES）指足以充分利用规模经济的最小产量，这一概念以平均成本曲线存在最小点为前提，即成本函数中含有规模不经济的成分、$AC$ 呈 U 形。本节 $AC(Q) = F/Q + c$ 的简化设定并不满足这一前提：$AC$ 对所有产量严格递减，并无最小点。但下降速率 $|AC'(Q)| = F/Q^2$ 随产量迅速衰减，产量足够大之后固定成本已经摊至可以忽略，继续扩张带来的单位成本节约微乎其微，产业组织实证因此把 $AC$ 落至接近边际成本 $c$ 的产量界定为最低有效规模，即图~\ref{fig:eos} 左图中的 $Q_{\text{MES}}$。\important{产量越过最低有效规模之后，大厂商相对小厂商的单位成本优势趋于消失，为多个厂商共存留下空间}。需求方规模经济则可能自我强化至\important{赢家通吃}（winner-take-all）：用户规模越大，产品价值越高，新用户越倾向于选择大网络，网络规模差距被正反馈不断放大，直至市场被单一网络占据（3.4 节）。需求方规模经济的这一性质决定了数字市场的竞争格局与反垄断难题（第 6 章），也是理解本章全部模型的线索。

\section{网络外部性的微观模型}

\subsection{基本模型：$u = \theta n - p$}

网络外部性建模的标准起点，是让消费者的效用直接依赖于网络规模。\mn{模型记忆：$u=\theta n - p$ 中 $\theta$ 是"网络偏好强度"，$n$ 是"预期/实际用户比例"，$p$ 是接入价格。推导题通常给出此效用函数，要求推导需求函数。}下面的模型是本章全部推导的基础，它依次回答三个问题：需求函数如何随预期规模变化、均衡如何由预期确定、为什么网络产品会呈现"要么大成、要么消亡"的极端结果。

\begin{derivation}{基本模型与临界规模}{}
考虑一个单一网络产品。消费者以偏好参数 $\theta$ 刻画对网络价值的异质性：类型为 $\theta$ 的消费者在网络规模为 $n$（标准化为用户比例，$n \in [0,1]$）时获得的效用为
\[
u(\theta) = \theta n - p,
\]
其中 $p$ 为接入价格。设 $\theta$ 在 $[0,1]$ 上均匀分布。

\textbf{第一步：需求函数。}给定预期网络规模 $n^e$，接入的边际消费者满足 $\theta n^e - p = 0$，即 $\hat\theta = p/n^e$。$\theta$ 高于 $\hat\theta$ 的消费者接入获得正效用，低于 $\hat\theta$ 的消费者不接入。由于 $\theta$ 在 $[0,1]$ 上均匀分布，接入用户占全体消费者的比例等于区间 $[\hat\theta, 1]$ 的长度，即 $1 - \hat\theta$。网络规模 $n$ 正是接入用户的比例，故需求为
\[
n = 1 - \hat\theta = 1 - \frac{p}{n^e}.
\]
给定价格 $p$，需求 $n$ 是预期规模 $n^e$ 的函数：
\[
n(n^e) = 1 - \frac{p}{n^e}.
\]
\textbf{第二步：均衡条件。}均衡要求预期自我实现（理性预期，rational expectations）：$n = n^e$。代入得
\[
n = 1 - \frac{p}{n} \;\Longleftrightarrow\; p = n(1-n).
\]
这就是网络外部性下的\important{反需求函数}：价格是用户规模的倒 U 形函数，在 $n = 1/2$ 处达到最大值 $p_{\max} = 1/4$。

\textbf{第三步：多重均衡。}对给定价格 $p < 1/4$，方程 $p = n(1-n)$ 有两个根（多重均衡，multiple equilibria）：
\[
n_1^* = \frac{1 - \sqrt{1-4p}}{2} < \frac12, \qquad
n_2^* = \frac{1 + \sqrt{1-4p}}{2} > \frac12.
\]
$n_1^*$ 为\important{临界规模}（critical mass，不稳定均衡），$n_2^*$ 为稳定均衡。均衡的稳定性由预期调整的动力学决定：设用户根据当前规模调整预期（$n_{t+1} = 1 - p/n_t$），则当 $n_t < n_1^*$ 时，实际需求低于预期，网络收缩趋向于零；当 $n_1^* < n_t < n_2^*$ 时，实际需求高于预期，网络扩张收敛于 $n_2^*$；当 $n_t > n_2^*$ 时，预期过度乐观，需求回落至 $n_2^*$。临界规模的经济含义：\important{网络必须越过 $n_1^*$ 才能启动正反馈并收敛于大网络均衡；低于临界规模，网络将因预期崩塌而消亡}。
\end{derivation}


\begin{figure}[htbp]
\centering
\begin{tikzpicture}[>=Stealth, scale=0.95]
\begin{axis}[
    width=10.5cm, height=6.5cm,
    xlabel={用户规模 $n$}, ylabel={价格 $p$},
    xmin=0, xmax=1.05, ymin=0, ymax=0.3,
    grid=major, axis lines=left,
    xtick={0,0.3,0.5,0.8,1}, ytick={0,0.1,0.2,0.25},
    legend pos=north east,
]
\addplot[color=main, very thick, domain=0:1, samples=100] {x*(1-x)};
\addlegendentry{$p = n(1-n)$ 反需求曲线}
\addplot[color=winered, dashed, thick] coordinates {(0,0.12) (1,0.12)};
\node[anchor=north, font=\small, color=winered] at (axis cs:0.5,0.16) {$p_0 < 1/4$};
\addplot[only marks, mark=*, color=second] coordinates {(0.1394,0.12)};
\node[anchor=north, font=\small, color=second] at (axis cs:0.1394,0.16) {$n_1^*$};
\addplot[only marks, mark=*, color=second] coordinates {(0.8606,0.12)};
\node[anchor=north, font=\small, color=second] at (axis cs:0.8606,0.16) {$n_2^*$};
\draw[->, color=second, thick] (axis cs:0.18,0.05) -- (axis cs:0.85,0.05);
\node[font=\small, anchor=south, color=second] at (axis cs:0.505,0.06) {正反馈扩张};
\end{axis}
\end{tikzpicture}
\caption{网络外部性下的倒 U 形需求与临界规模}
\label{fig:network}
\end{figure}

图~\ref{fig:network} 直观呈现了多重均衡的结构。价格 $p_0$ 水平线与倒 U 形反需求曲线的两个交点中，左侧交点是不稳定的临界规模，右侧交点是稳定均衡；两交点之间的用户规模区间，正反馈将推动网络扩张至稳定均衡，临界规模之下的网络则收缩至零，稳定均衡右侧的过度乐观预期同样会回落至右侧均衡点$n_2^*$。这一图形结构解释了现实中网络产品"要么成功要么消亡"的极端命运：网络外部性使中间状态不稳定。

\subsection{预期协调：Katz–Shapiro 模型}

基本模型中，多重均衡的存在提出了预期协调问题：既然多个均衡都可能实现，最终哪一个会被选中？Katz 与 Shapiro 的贡献在于提出\textbf{预期满足均衡}（fulfilled-expectations equilibrium）的概念：预期在均衡选择中起决定作用——消费者预期的网络规模本身就是均衡的组成部分，\important{乐观预期自我实现为大规模均衡，悲观预期自我实现为零规模均衡}。

\begin{derivation}{预期满足均衡}{}
沿用 $u = \theta n - p$ 的设定，预期满足条件为 $n = n(n^e) = 1 - p/n^e$ 与 $n^e = n$ 同时成立，即不动点方程
\[
n = 1 - \frac{p}{n}.
\]
该方程的根即全部理性预期均衡。其经济含义是：\important{价格相同的同一产品，可能因预期差异而得到完全不同的市场结果}——若预期 $n^e < n_1^*$，实际需求 $n(n^e)$ 低于预期，预期下调、需求进一步收缩，网络消亡；若预期位于 $n_1^*$ 与 $n_2^*$ 之间，实际需求高于预期，预期上调、需求扩张，直至 $n_2^*$；若预期 $n^e > n_2^*$，实际需求低于预期，回落至 $n_2^*$。预期协调问题的政策含义深远：厂商有强烈动机操纵预期——预发布（announcement）、补贴首期用户、兼容策略（3.3 节），本质上都是预期管理。平台"烧钱换增长"的合理性正在于此：补贴的目的在于把消费者预期推过临界规模，当期利润的让渡是为此付出的代价。
\end{derivation}

网络外部性使新产品面临\textbf{冷启动问题}（cold start problem）：网络规模为零时产品对任何人都没有价值，无人愿意成为第一个用户，"先有鸡还是先有蛋"的困局由此而来。\mn{预期管理三件套：预发布、补贴首期用户、兼容策略。"烧钱换增长"的实质＝把消费者预期推过临界规模。}跨越临界规模的常见策略均围绕预期与初始用户展开：补贴与免费获客是最直接的手段，微信支付借 2015 年央视春晚的"摇一摇"红包活动在数日内完成了大规模的用户支付绑定，滴滴与快的的补贴大战同样是付费购买初始网络；预装与渠道合作（手机预装、系统捆绑）可以一次性获得用户基数；预发布造势则用未来预期撬动当期接入，即提前宣布产品、开放预约，让消费者相信网络终将成形；此外，新产品还可以兼容既有大网络（3.3 节），把对方的存量用户变成自己的初始网络。这些策略的本质都是预期管理：把消费者预期推过临界规模，正反馈随后接管（第 4 章将给出补贴的最优定价分析）。

\subsection{路径依赖与锁定：QWERTY 的教训}

预期决定均衡，而预期又受历史路径影响。为什么消费者会预期某个网络胜出？答案往往要到已经走过的历史中去找：今天的用户规模是过去无数个体决策的累积结果，历史中的微小偶然事件可能把网络引向不同的轨道。这就是路径依赖问题。

\begin{definition}[路径依赖（path dependence）]
路径依赖是指微小历史事件通过自我强化机制将长期结果锁定（lock-in）、且被锁定结果未必最优的现象：初始状态的微小差异经由正反馈被放大为长期的质的差异。
\end{definition}

QWERTY 键盘是路径依赖的标准案例。\mn{案例锚点：1985 年 David 的经典论文以 QWERTY 论证路径依赖"锁定"；Liebowitz–Margolis 反驳 QWERTY 与 Dvorak 的效率差异缺乏实证。}现代键盘布局在打字机时代形成，其设计初衷是把常用字母组合在机械上彼此分散、减少打字机连杆卡键，客观上牺牲了部分打字速度。电子键盘早已不存在卡键问题，理论上更高效的布局（如 Dvorak）却始终未能取代 QWERTY，因为所有打字者的技能、教材、培训体系都建立在 QWERTY 之上，个体转换布局的成本远高于效率收益。路径依赖的形成机制在于\textbf{自我强化}：一种选择的使用者越多，其配套资源（技能、教材、兼容设备）越多，新使用者越倾向于选择它。网络外部性放大历史偶然事件的影响，使"好的"技术不一定胜出。

路径依赖之所以能够长期维持，关键在于用户脱离既有网络要付出代价。这一代价在理论上被概括为转换成本，它与网络规模差距共同构成在位网络的护城河。

\begin{definition}[转换成本（switching cost）]
转换成本是用户从一个产品或网络迁移到另一个产品或网络时必须一次性承担的代价，包括学习新系统的时间成本、数据与社会关系的迁移成本、配套设备的更换成本，以及迁移期间服务中断的损失。
\end{definition}

转换成本存在时，挑战者要撬动在位网络的用户，就必须在价格上让出足够的空间。让出多少才够，取决于转换成本本身，也取决于两个网络的规模差距。下面的推导把网络外部性显式纳入效用函数，求出挑战者的定价上限、必须让出的价格差，以及价格手段彻底失效的临界条件。

\begin{derivation}{网络外部性下的迁移壁垒与锁定}{}
设消费者当前使用在位网络 A，正在考虑迁移到挑战者网络 B。两个网络的规模分别为 $N_A$ 与 $N_B$（标准化为用户比例），且在位者规模占优，$N_A > N_B$。消费者的网络偏好强度为 $\theta > 0$，两个产品的单机价值（与网络规模无关的内在价值）分别为 $v_A$ 与 $v_B$，接入价格分别为 $p_A$ 与 $p_B$。消费者使用网络 $i$ 的效用为
\[
u_i = \theta N_i + v_i - p_i, \qquad i \in \{A, B\},
\]
迁移时一次性支付转换成本 $s > 0$。

\textbf{第一步：转换条件。}消费者迁移的前提，是迁移后的净效用不低于留在原网络的效用。迁移所得为 $u_B$，但要扣去转换成本 $s$，故转换条件为
\[
u_B - s \geq u_A.
\]
把效用函数代入，将网络价值、单机价值与价格三项全部展开，得
\[
\theta N_B + v_B - p_B - s \geq \theta N_A + v_A - p_A.
\]

\textbf{第二步：挑战者的定价上限。}上式对挑战者构成一个定价约束：只有把价格定得足够低，不等式才能成立。将含 $p_B$ 的项移到一侧、其余项移到另一侧，解出
\[
p_B \leq p_A - \theta(N_A - N_B) + (v_B - v_A) - s.
\]
不等式右端即挑战者能够吸引用户迁移的\important{最高定价} $p_B^{\max}$。它以在位者价格 $p_A$ 为基准，被三项修正：网络规模劣势 $\theta(N_A - N_B)$ 与转换成本 $s$ 各自压低这一上限，单机价值优势 $v_B - v_A$ 则抬高它。

\textbf{第三步：临界价格差。}定义挑战者必须相对在位者让出的价格差
\[
\Delta \equiv p_A - p_B^{\max} = \theta(N_A - N_B) + s - (v_B - v_A),
\]
$\Delta$ 度量迁移壁垒的货币价值：B 必须把价格定在 $p_A - \Delta$ 及以下才能撬动用户。$\Delta > 0$ 意味着 $p_B^{\max} < p_A$，即挑战者不降价就留不住迁移意愿，故在双方定价相同（$p_B = p_A$）的情形下，$\Delta > 0$ 恰是在位者维持锁定的条件。在此基础上再令两网络规模相同（$N_A = N_B$），网络价值项消失，$\Delta = s - (v_B - v_A)$；由于此时 $u_B - u_A = (v_B - v_A) - (p_B - p_A) = v_B - v_A$，锁定条件退化为
\[
v_B - v_A < s \quad\Longleftrightarrow\quad u_B - u_A < s,
\]
即\important{挑战者的效用优势小于转换成本时，在位网络维持锁定}。网络规模存在差距时，$\theta(N_A - N_B) > 0$ 使 $\Delta$ 在纯转换成本之上再加一层，迁移壁垒被网络外部性系统性抬高。

\textbf{第四步：绝对锁定条件。}价格不能为负，挑战者的可行定价须满足 $p_B \geq 0$，因此 B 能够撬动用户的前提是 $p_B^{\max} = p_A - \Delta \geq 0$。反过来，当
\[
\Delta > p_A, \quad \text{即} \quad \theta(N_A - N_B) + s - (v_B - v_A) > p_A
\]
时 $p_B^{\max} < 0$，意味着即使 B 定价为零，转换条件仍不成立，\important{在位网络的锁定无法被价格手段打破}。在位者本身免费（$p_A = 0$）时这一条件最容易满足，只要 $\theta(N_A - N_B) + s > v_B - v_A$ 即告成立。

\textbf{第五步：比较静态。}对 $\Delta$ 分别求偏导，考察三个参数各自的作用方向：
\[
\frac{\partial \Delta}{\partial s} = 1 > 0, \qquad
\frac{\partial \Delta}{\partial (N_A - N_B)} = \theta > 0, \qquad
\frac{\partial \Delta}{\partial (v_B - v_A)} = -1 < 0.
\]
转换成本与网络规模差距同向抬高迁移壁垒，单机价值优势反向削弱它。令 $\Delta$ 保持不变并对前两个参数求全微分，$\theta \,\mathrm{d}(N_A - N_B) + \mathrm{d}s = 0$，得二者的替代率
\[
\frac{\mathrm{d}s}{\mathrm{d}(N_A - N_B)} \bigg|_{\Delta} = -\theta,
\]
即\important{网络规模差距每扩大一单位，在位者可以少依赖 $\theta$ 单位的转换成本来维持同等强度的锁定}。
\end{derivation}

推导的结论可以概括为一句话：\important{网络市场上的迁移壁垒等于转换成本加网络价值差}。纯转换成本模型只解释了 $s$ 这一项，$\theta(N_A - N_B)$ 则说明，即使挑战者的产品本身更好、转换成本也不算高，只要在位网络的规模优势足够大，用户仍然不会迁移。\mn{锁定三问：转换成本多高（$s$）？网络规模差多大（$\theta(N_A-N_B)$）？产品好多少（$v_B-v_A$）？三者决定挑战者要让出的价格差 $\Delta$。}$\Delta$ 的表达式同时给出了挑战者可用的两类手段：一类是在价格上让利，把 $p_B$ 压到 $p_A - \Delta$ 以下；另一类是改变 $\Delta$ 本身，即提高单机价值 $v_B$、降低转换成本 $s$，或者压缩网络规模差距 $N_A - N_B$。第四步的条件说明，当 $\Delta$ 超过在位者价格 $p_A$ 时第一类手段彻底失效，挑战者只剩第二类可用。兼容策略正是压缩 $N_A - N_B$ 的手段：接入在位者的网络之后，对方的存量用户同时成为自己的有效网络规模（3.3 节详述）。

中国互联网的一组对照可以印证这一机制。2019 年 1 月，多闪、聊天宝、马桶 MT 三款社交产品在同一天发布，功能定位各有差异且全部免费，但均未能撼动微信的地位。用推导的语言说，微信本身免费（$p_A = 0$），挑战者在价格上无从让利，而社交关系链既构成巨大的网络价值差 $\theta(N_A - N_B)$，又意味着高昂的关系重建成本 $s$，$\Delta > p_A$ 的绝对锁定条件成立。反向的例子来自监管干预：工业和信息化部自 2019 年 11 月起在全国推行携号转网，把更换运营商所需的号码迁移成本从制度上压到接近于零，正是通过降低 $s$ 来削弱在位运营商的锁定。这提示了数字市场竞争政策的一个着力点：降低用户迁移成本可以在不改变市场结构的前提下增强竞争压力（第 6 章）。

\section{兼容性与标准化}

\subsection{兼容性决策}

面对网络外部性，厂商的核心战略选择是\textbf{兼容}（compatibility）还是\textbf{不兼容}。兼容意味着产品接入共同的网络，用户的网络价值取决于全网规模；不兼容则各产品构成独立网络，用户价值取决于本产品规模。兼容性决策的本质是市场规模的交换：兼容扩张了有效网络规模（需求方规模经济做大），但抹平了产品间网络规模的差异（削弱差异化优势）。

\begin{derivation}{兼容性决策}{}
考虑两家厂商，产品完全同质，边际成本为 $c$，唯一差异是网络规模。设不兼容时各厂商的网络规模分别为 $N_1 > N_2$（历史原因形成的不对称），兼容时有效网络规模为 $N = N_1 + N_2$。消费者的效用为
\[
u_i = \theta N_i + v - p_i,
\]
其中 $v$ 为产品内在价值，$\theta N_i$ 为网络价值。不兼容时，大网络厂商可以凭借 $N_1 > N_2$ 收取溢价：在 $v$ 相同的条件下，消费者为接入网络 1 愿意多支付 $\theta(N_1 - N_2)$。大网络厂商由此获得的利润优势约为
\[
\Delta\pi_{\text{不兼容}} \approx \theta(N_1 - N_2) \cdot q_1 > 0,
\]
其中 $q_1$ 为大网络厂商的销量（价格溢价乘以销量）。兼容后，两厂商网络价值相同，竞争退化为纯粹的同质产品价格竞争：$p_1 = p_2 = c$，网络溢价消失；兼容的收益在于总网络规模 $N$ 扩大使全体消费者的支付意愿上升、需求整体扩张（这正是需求方规模经济的含义），代价则是网络规模差异这一竞争优势被抹平。因此：\important{网络规模占优的厂商倾向于不兼容（保住溢价），网络规模落后的厂商倾向于兼容（抹平差距）}；当产品高度差异化（各自的网络价值针对不同用户群）时，兼容更可能被双方接受。这一不对称解释了标准竞争中大厂与小厂的典型立场：\important{大厂要封闭，小厂要开放}。
\end{derivation}

兼容性决策的经典中国案例是 WPS 与微软 Office 的竞争：WPS 长期以深度兼容 Office 文档格式（.doc/.docx）为生存策略——小厂商主动兼容大网络、抹平网络规模差距，这正是"小厂要开放"的逻辑；\mn{兼容决策口诀：大厂要封闭（保溢价），小厂要开放（抹差距）；产品高度差异化时兼容更易达成。}微软则通过格式的封闭与不断迭代维持其事实标准地位，体现了"大厂要封闭"的立场。反向的例子是游戏主机：索尼、微软、任天堂的主机长期互不兼容、游戏库彼此独占，因为三家的用户群与游戏生态高度差异化，封闭保住了各自的差异化优势，而兼容带来的总网络收益有限。

\subsection{标准竞争}

兼容性决策在产业层面表现为标准竞争：厂商争夺的不仅是今天的用户，更是让自家的技术方案成为整个行业共同遵循的标准——标准一旦胜出，其网络效应将带来持久的市场地位。

\begin{definition}[标准（standard）]
标准是保证产品互操作性（interoperability）的技术规范，分为\textbf{正式标准}（de jure standard，标准化组织制定，如 5G 标准）与\textbf{事实标准}（de facto standard，市场自发形成，如 Windows 的操作系统地位）。
\end{definition}

标准竞争的特殊性在于其"赢家通吃"性质与网络外部性叠加：标准一旦确立，其使用者网络构成巨大的进入壁垒，标准竞争的胜者获得长期的市场主导地位。

标准必要专利（standard-essential patent，SEP）是标准竞争中的关键制度安排：当某一专利技术被纳入标准、任何遵循标准的产品都必须使用时，该专利成为标准必要专利。为防止专利持有人借标准的锁定效应"敲竹杠"（hold-up），标准化组织要求 SEP 持有人承诺以\textbf{公平、合理、无歧视}（fair, reasonable and non-discriminatory，FRAND）的条件许可专利。\mn{SEP/FRAND 是数字治理与反垄断的热点：SEP 持有人拒绝许可或索取过高许可费构成滥用市场支配地位（第 6 章）。华为与诺基亚、苹果与高通等纠纷均围绕 SEP 许可费。}FRAND 承诺的经济学含义：在标准已经形成锁定之后，用事前承诺约束事后的垄断定价，这是对"锁定后敲竹杠"问题的制度性防御。

\begin{derivation}{锁定后敲竹杠}{}
设标准确立之前，专利持有人与实施者的谈判力量对称，事前竞争的许可费为 $r_c$（接近专利技术对产品的边际贡献）。标准确立后，实施者为脱离该专利必须承担转换成本 $S$（重新设计、产线调整、市场损失）。此时专利持有人能够索取的最高许可费满足实施者"留用与退出无差异"：
\[
r_{\max} = r_{\text{替代方案}} + S,
\]
其中 $r_{\text{替代方案}}$ 为采用替代技术的成本。事前许可费与事后上限之差
\[
\Delta r = r_{\max} - r_c \approx S + (r_{\text{替代方案}} - r_c)
\]
即\textbf{锁定租金}：专利持有人事后索取的许可费可以系统性地超过事前竞争水平，超出部分由实施者的转换成本 $S$ 支撑。FRAND 承诺的制度功能正在于此：它把事后定价约束回到事前竞争的参照系（"公平合理"的标准解释即为"事前竞争水平下双方自愿达成的价格"），从而阻止锁定租金的提取。这一框架同时解释了 SEP 诉讼的标准争议：许可费纠纷的实质，是"事前竞争价格"在标准确立后已不可观察，只能通过可比许可、最小可销售单元等反事实方法近似还原。
\end{derivation}

\section{网络外部性与市场结构}

\subsection{赢家通吃的条件}

网络外部性使市场可能出现"赢家通吃"结构：单一厂商占据绝大部分市场。但并非所有网络市场都会走向赢家通吃，其成立需要若干条件同时满足。

\textbf{条件一：强网络外部性。}网络价值对用户规模的弹性足够大，大网络对小网络的优势无法被价格或差异化抵消。

\textbf{条件二：供给方规模经济。}固定成本高、边际成本低的成本结构使大厂商具有持续的成本优势：需求方与供给方规模经济叠加，赢家通吃的力量最强（搜索引擎、操作系统兼具两者）。

\textbf{条件三：用户偏好同质。}若用户对产品特征的偏好高度异质，垂直细分市场可以为小网络提供庇护：小众社交应用凭借差异化生存，说明偏好异质性是对赢家通吃的天然制约。

\textbf{条件四：多归属成本高。}若用户可以低成本地同时使用多个网络（多归属，multi-homing），小网络不会因规模劣势被抛弃，市场趋于多平台共存（第 4 章详述）；只有当多归属成本高、用户被迫单归属时，网络规模优势才会转化为赢家通吃。\mn{赢家通吃四条件：强网络外部性、供给方规模经济、用户偏好同质、多归属成本高。反向检验：多归属成本低 → 多平台共存（第 4 章）。}国内市场的对照同样鲜明：即时通讯领域微信一家独大（强网络外部性、高转换成本、事实上的单归属），而外卖、网约车领域美团与饿了么、滴滴与高德长期多平台共存，后者的多归属成本几乎为零，用户同时安装多个 App 毫无障碍。

\subsection{网络外部性的福利含义}

赢家通吃刻画的是网络外部性对竞争结构的影响；效率维度的评价则要回答另一个问题：网络外部性存在时，市场均衡的网络规模是否是社会最优的？直觉提示答案可能是否定的：每个用户的接入都使所有存量用户受益，而接入决策只考虑私人收益，这一正外部性未被内部化。基本模型可以精确刻画这一偏差。

\begin{derivation}{网络外部性的福利分析}{}
\textbf{网络外部性下的市场失灵：福利分析。}设反需求函数 $p = n(1-n)$（基本模型推导所得）。给定接入价格 $p_0$，接入的边际用户（类型 $\hat\theta = 1-n$）满足私人支付意愿等于价格：
\[
\hat\theta n = n(1-n) = p_0,
\]
市场均衡网络规模即由此确定（多重均衡中取稳定均衡 $n_2^*$）。

\textbf{消费者总剩余的度量。}反需求曲线 $p(x) = x(1-x)$ 给出第 $x$ 个用户（边际用户）的支付意愿：其类型为 $1-x$，网络规模为 $x$ 时网络价值为 $(1-x)\,x = x(1-x)$。所有用户支付同一价格 $p_0$，故网络规模 $n$ 处的消费者总剩余等于"支付意愿减去实付价格"沿用户累积，即
\[
CS(n) = \int_0^n \bigl(x(1-x) - p_0\bigr) \, dx.
\]

\textbf{社会最优。}社会最优要求\textbf{边际社会收益}等于价格。网络的外部性意味着最后一个用户的接入同时提升了所有存量用户的效用，但这一外部收益在接入决策中未被考虑。第 $n$ 个用户接入带来的边际社会收益等于其私人支付意愿 $n(1-n)$ 加上他给存量用户带来的价值增量：每个存量用户 $i$ 的效用为 $\theta_i n - p_0$，一个新用户加入使每个存量用户的效用提升 $\theta_i$；存量用户构成类型区间 $[1-n, 1]$，故总外部收益为
\[
\int_{1-n}^{1} \theta \, d\theta = \frac{1 - (1-n)^2}{2} = n - \frac{n^2}{2}.
\]
于是边际社会收益
\[
MSB(n) = n(1-n) + n - \frac{n^2}{2} = 2n - \frac{3}{2} n^2,
\]
在任何 $n \in (0,1)$ 处均高于私人支付意愿 $n(1-n)$（二者之差为 $n - \tfrac12 n^2 > 0$），故社会最优网络规模严格大于市场均衡规模；特别地，$MSB(1) = 1/2 > p_{\max} = 1/4 \geq p_0$，边际社会收益在网络规模上界处仍高于价格，社会最优为全市场覆盖（$n = 1$），与市场均衡的 $n_2^* < 1$ 形成鲜明对照。结论：\important{网络外部性为正时，市场均衡的网络规模小于社会最优规模}，这是庇古外部性理论的直接推论：正外部性导致供给不足。矫正方向包括补贴接入（降低 $p$）、政府投资网络基础设施（公共品供给），或对网络运营者给予排他性激励（专利制度），使外部性内部化。
\end{derivation}

\subsection{网络垄断与传统垄断}

赢家通吃下的垄断与传统产业组织意义上的垄断成因不同，后果也不同，二者的区分是数字反垄断的理论入口。成因上，传统垄断来自\textbf{供给方}力量：规模经济、专利壁垒、行政准入，市场份额来自成本优势；网络垄断来自\textbf{需求方}力量：网络外部性使大网络对小网络具有不可替代的价值优势，市场份额来自用户的彼此依赖。后果上，传统垄断以限产提价损害消费者（垄断价格高于竞争价格、产量低于竞争产量），而网络垄断的福利效应是混合的：垄断平台的大网络本身就是用户价值的来源，拆解大网络反而降低网络效用，反垄断的着力点因此落在行为规制（排他、拒绝接入、自我优待）上，"反大即反垄断"的传统逻辑与结构拆解手段在这里都不适用。\mn{网络垄断 vs 传统垄断一句话对照：成因看需求方 vs 供给方，后果看网络价值 vs 限产提价，规制看行为 vs 结构。第 6 章反垄断的分析框架以此区分展开。}这一区分对应着第 6 章的核心命题：网络效应同时是市场势力的来源与消费者福利的来源，规制工具的选择必须在这两者之间权衡。

\subsection{本章的政策逻辑}

网络外部性的政策含义构成数字经济监管的理论起点：其一，市场可能收敛于劣等的网络或标准（路径依赖），政府可通过标准化政策、互操作性要求矫正；其二，临界规模的存在意味着市场进入困难，先发优势可能固化垄断（第 6 章反垄断）；其三，正外部性意味着私人供给不足，数字基础设施的公共投资具有效率依据。这些结论将在第 4 章（平台）、第 6 章（反垄断）与第 12 章（治理）中被反复调用。

\begin{chapterbox}
\textbf{本章考点清单}

\begin{enumerate}
\item 网络外部性的定义、直接与间接之分、正负之分、与网络效应的辨析 \important{【专硕·核心】}
\item 需求方规模经济与供给方规模经济的区分；最低有效规模与供给方规模优势的衰减 \important{【专硕·核心】}
\item $u = \theta n - p$ 模型：需求函数 $n = 1 - p/n^e$、反需求函数 $p = n(1-n)$ 的推导 \important{【专硕·核心】}
\item 临界规模 $n_1^* = (1-\sqrt{1-4p})/2$ 的推导与均衡稳定性分析 \important{【专硕·核心】}
\item Katz–Shapiro 预期满足均衡：预期操纵与补贴获客的理论依据；冷启动问题与跨越临界规模的策略 \important{【专硕·核心】}
\item 路径依赖与 QWERTY 案例；迁移壁垒 $\Delta = \theta(N_A - N_B) + s - (v_B - v_A)$ 的推导、锁定条件 $u_B - u_A < s$ 与绝对锁定条件 $\Delta > p_A$ \important{【专硕·核心】}
\item 兼容性决策：大厂商封闭、小厂商开放的利润比较 \important{【专硕·核心】}
\item 标准、事实标准、SEP 与 FRAND 承诺 \important{【专硕·扩展】}
\item 赢家通吃的四个条件 \important{【专硕·核心】}
\item 网络外部性的福利含义：正外部性导致网络规模低于社会最优 \important{【专硕·核心】【学硕·热点】}
\item 网络垄断 vs 传统垄断：需求方 vs 供给方成因、网络价值 vs 限产提价后果、行为规制 vs 结构拆解 \important{【专硕·核心】}
\end{enumerate}
\end{chapterbox}

本章的模型揭示了一个根本矛盾：网络外部性使市场力量自我强化，而这一力量既可以成就卓越的平台，也可以固化劣等的均衡。把握这一矛盾——价值如何随规模加速增长、均衡如何被预期锁定——是理解后续平台竞争与监管问题的共同前提。
